% ETC: exact constants in the original CCM degree-adjusting test family.
\section{The original trivial-correspondence family with every Haar constant}
\label{sec:cc-exact-trivial-correspondence}

\paragraph{Human source and reading scope.}
Alain Connes, Caterina Consani, and Matilde Marcolli,
\href{https://arxiv.org/abs/math/0703392}{\emph{The Weil proof and the
geometry of the adeles class space}}, original author file
\nolinkurl{Yuri70v2.tex}, subsection \emph{Fubini's theorem and the trivial
correspondences}, especially Lemma \texttt{degmodifyV}, supply the
polynomial Gaussian used below. The original file has SHA-256
\nolinkurl{b1f694cb934dd4fef829287f1a55dd8de21c67a481e552453c2f7e8fe0ae6fb5}.
The complete subsection, source lines 2313--2424, was read, together
with the definitions \texttt{fsharpinv}, \texttt{degcorrZf}, and
\texttt{codegcorrZf}. This calculation fixes the measures and sums over
all nonzero rational numbers; it replaces the source's unspecified
constant by a proved exact value. It does not replace the original
zeta function with a completed function, and makes no RH inference.

\subsection{Original adelic test, rational restrictions, and additive moments}

Use additive Haar measure $dx$ at infinity and
$\operatorname{vol}(\mathbb Z_p)=1$ at every finite prime. On
$C_{\mathbb Q}=\widehat{\mathbb Z}^{\times}\times\mathbb R_{>0}$
use Haar measure $du\,dy/y$, with
$\operatorname{vol}(\widehat{\mathbb Z}^{\times})=1$. Retain
\begin{equation}\tag{ETC1}
 \eta(x)=\pi x^2\left(\pi x^2-\frac32\right)e^{-\pi x^2}
        =\left(\pi^2x^4-\frac{3\pi}{2}x^2\right)e^{-\pi x^2},
 \quad k\in\bigoplus_p\mathbb Z,
 \quad N(k)=\prod_pp^{k_p},\quad a>0,
\end{equation}
and the actual tensor
\begin{equation}\tag{ETC2}
 \xi_{k,a}(x)=\eta(x_\infty/a)
                  \prod_p1_{p^{k_p}\mathbb Z_p}(x_p),
 \qquad b=\frac a{N(k)}>0.
\end{equation}
The symbol $b$ records this specified ratio; it does not change either
of the original parameters or the local factors. The finite product
of nonstandard factors defines a Bruhat--Schwartz function on the
original adeles.

The Gaussian integral and its two differentiated forms are
\begin{equation}\tag{ETC3}
 \int_{\mathbb R}e^{-\pi x^2}\,dx=1,\quad
 \int_{\mathbb R}x^2e^{-\pi x^2}\,dx=\frac1{2\pi},\quad
 \int_{\mathbb R}x^4e^{-\pi x^2}\,dx=\frac3{4\pi^2}.
\end{equation}
For example, square the first integral and use polar coordinates to
obtain one; differentiating
$\int e^{-\pi t x^2}dx=t^{-1/2}$ once and twice at $t=1$
gives the other two. Polynomial Gaussian bounds justify both
differentiations. Consequently
\begin{equation}\tag{ETC4}
 \eta(0)=0,\quad \int_{\mathbb R}\eta=\frac34-\frac34=0,
 \quad \xi_{k,a}(0)=0,\quad
 \int_{\mathbb A_{\mathbb Q}}\xi_{k,a}
 =a\prod_pp^{-k_p}\left(\frac34-\frac34\right)=0.
\end{equation}
Thus the two additive zero-moment conditions hold on the original
test, with both summands and every volume retained.

Define the reduction by the original full rational sum:
\begin{equation}\tag{ETC5}
 f_{k,a}(u,y)=(\mathcal E\xi_{k,a})(u,y)
 =\sum_{q\in\mathbb Q^\times}\xi_{k,a}(qy,(qu_p)_p).
\end{equation}
Since $u_p$ is a unit, its finite characteristic factors permit exactly
$v_p(q)\ge k_p$ at every prime. This says that $q/N(k)$ is a rational
number integral at every finite prime, hence an integer: a reduced
denominator greater than one would give a negative valuation at one
of its prime divisors. Conversely every nonzero integer satisfies
these conditions. The map $n\mapsto N(k)n$ is therefore a bijection
onto the permitted rational numbers. Since $\eta$ is even,
\begin{equation}\tag{ETC6}
 f_{k,a}(u,y)=\sum_{n\in\mathbb Z\setminus\{0\}}
             \eta\left(\frac{N(k)ny}{a}\right)
 =2\sum_{n=1}^{\infty}\eta(ny/b)=:f_b(y).
\end{equation}
The factor two is from the two rational signs. This is an exact map
of the specified family; it shows explicitly that this scalar
reduction depends on $k,a$ only through $a/N(k)$. It does not recover
the individual finite-place factors from that ratio.

\subsection{Fourier calculation, Poisson identity, and convergence}

Use the real Fourier convention
$\widehat h(v)=\int h(x)e^{-2\pi ixv}dx$ and compatible conductor
$\mathbb Z_p$ characters at finite places. Let $g(x)=e^{-\pi x^2}$.
The Fourier integral has derivative $-2\pi v\widehat g(v)$:
differentiate under the integral and integrate $g'=-2\pi xg$ by
parts. Its value at zero is one by ETC3, so $\widehat g(v)=g(v)$.
Differentiating this identity gives the complete two polynomial
transforms
\begin{equation}\tag{ETC7}
 \begin{split}
 \widehat{x^2g}(v)&=\left(\frac1{2\pi}-v^2\right)g(v),\\
 \widehat{x^4g}(v)&=
       \left(v^4-\frac3\pi v^2+\frac3{4\pi^2}\right)g(v),\\
 \widehat\eta(v)&=\left[
 \pi^2\left(v^4-\frac3\pi v^2+\frac3{4\pi^2}\right)
 -\frac{3\pi}{2}\left(\frac1{2\pi}-v^2\right)\right]g(v)
 =\eta(v).
 \end{split}
\end{equation}
These differentiations and integrations by parts have vanishing
boundary terms because every derivative is a polynomial times a
Gaussian. At a finite prime, character orthogonality gives
$\widehat{1_{p^{k_p}\mathbb Z_p}}
 =p^{-k_p}1_{p^{-k_p}\mathbb Z_p}$: the transform equals the
original ball volume on its annihilator and is zero outside it.
The real dilation contributes $a$. Hence
\begin{equation}\tag{ETC8}
 \widehat\xi_{k,a}
 =a\prod_pp^{-k_p}\,\xi_{-k,1/a}
 =\frac a{N(k)}\,\xi_{-k,1/a}.
\end{equation}
Applying Fourier twice yields the original even test: the second
coefficient is $(1/a)/N(-k)=N(k)/a$, so the two displayed coefficients
multiply to one.

For the required Poisson identity no rearrangement of a divergent
integral is used. For $r>0$, periodize $\eta(r(x+n))$ over
$n\in\mathbb Z$. Every differentiated series converges uniformly
on $x\in[0,1]$ by Schwartz bounds. Its $m$th Fourier coefficient
is $r^{-1}\widehat\eta(m/r)$, by the absolutely convergent
substitution over all intervals. These coefficients decrease faster
than any power, so their Fourier series converges absolutely to the
periodization. Evaluating at zero and using ETC7 gives
\begin{equation}\tag{ETC9}
 \sum_{n\in\mathbb Z}\eta(nr)
 =\frac1r\sum_{m\in\mathbb Z}\eta(m/r),\qquad
 f_b(y)=\frac b y f_{1/b}(1/y).
\end{equation}
The zero summand is zero on both sides because $\eta(0)=0$.
The same argument with $g$ proves its Gaussian Poisson formula.

For $y\ge1$ the series in ETC6 and each of its derivatives are
bounded by a polynomial in $y$ times $e^{-\pi y^2/b^2}$.
Indeed factor this last exponential from each term; the remaining
sum is bounded by a convergent series of a polynomial in $n$ times
$e^{-\pi(n^2-1)/b^2}$. ETC9 therefore bounds $f_b$ near zero
by a constant times $y^{-5}e^{-\pi b^2/y^2}$; its differentiated
versions have the same exponential with possibly different finite
powers of $y^{-1}$. Thus this actual reduced test and all its
multiplicative derivatives decrease faster than every power at both
ends. In particular every Mellin integral below converges absolutely
for every $s$, locally uniformly with all $s$ derivatives.

\subsection{Original-zeta Mellin formula with both Gaussian summands}

With the specified class-group measure put
\begin{equation}\tag{ETC10}
 M_b(s)=\int_{\widehat{\mathbb Z}^{\times}}du
            \int_0^\infty f_b(y)y^s\frac{dy}{y}
       =\int_0^\infty f_b(y)y^{s-1}dy.
\end{equation}
The preceding bounds prove that $M_b$ is entire. Initially for
$\operatorname{Re}s>1$, absolute Fubini is valid on ETC6, since
\[
 2\sum_{n\ge1}\int_0^\infty
    |\eta(ny/b)|y^{\operatorname{Re}s-1}dy
 =2b^{\operatorname{Re}s}\sum_{n\ge1}n^{-\operatorname{Re}s}
        \int_0^\infty|\eta(x)|x^{\operatorname{Re}s-1}dx<\infty.
\]
The local integral converges already for $\operatorname{Re}s>-2$,
as $\eta(x)=O(x^2)$ at zero. Substituting $v=\pi x^2$ into
each of the two original polynomial-Gaussian terms proves
\begin{equation}\tag{ETC11}
 \begin{split}
 \int_0^\infty\eta(x)x^{s-1}dx
 &=\frac{\pi^2}{2\pi^{(s+4)/2}}
       \Gamma\left(\frac{s+4}{2}\right)
   -\frac{3\pi}{4\pi^{(s+2)/2}}
       \Gamma\left(\frac{s+2}{2}\right),\\
 M_b(s)
 &=2\left(\frac a{N(k)}\right)^s\zeta(s)
   \left[
     \frac{\pi^2}{2\pi^{(s+4)/2}}
       \Gamma\left(\frac{s+4}{2}\right)
    -\frac{3\pi}{4\pi^{(s+2)/2}}
       \Gamma\left(\frac{s+2}{2}\right)
   \right].
 \end{split}
\end{equation}
The working expression keeps the original $\zeta(s)$, both Gamma
summands, the two rational signs, and all finite and real factors.
It continues meromorphically by the original zeta continuation
ZTC10; equality with the entire integral ETC10 proves that its
singularities are removable as a whole, without asserting that its
separate summands are entire.

For a precise comparison with the source's unspecified constant,
the two-step Gamma recurrence gives the identity
\begin{equation}\tag{ETC12}
 2\left[
     \frac{\pi^2}{2\pi^{(s+4)/2}}\Gamma\left(\frac{s+4}{2}\right)
    -\frac{3\pi}{4\pi^{(s+2)/2}}\Gamma\left(\frac{s+2}{2}\right)
 \right]
 =\frac{s(s-1)}4\pi^{-s/2}\Gamma(s/2).
\end{equation}
This is an explicit equality of the retained multiplier, not a
replacement arithmetic function. For real $u$ it yields the exact
source comparison
\begin{equation}\tag{ETC13}
 M_b(iu)=-\frac{b^{iu}}4\,u(u+i)
                   \pi^{-iu/2}\Gamma(iu/2)\zeta(iu).
\end{equation}
At removable points its value is the entire continuation specified
by ETC10. The factor $-1/4$ is for the full rational sum in ETC5.
Using only the positive integers would instead give $-1/8$ and
would not be the reduction defined in ETC5.

\subsection{Both endpoint values without an illicit interchange}

Retain the complete Gaussian theta sum, including its zero term:
\begin{equation}\tag{ETC14}
 \Theta_b(y)=\sum_{n\in\mathbb Z}e^{-\pi(ny/b)^2},\quad
 f_b(y)=\frac14\bigl(y^2\Theta_b''(y)+2y\Theta_b'(y)\bigr)
       =\frac14\bigl(y^2\Theta_b'(y)\bigr)'.
\end{equation}
For each summand this follows from
$x^2g''(x)+2xg'(x)=4\eta(x)$; the differentiated sums converge
locally uniformly. Gaussian Poisson summation as proved after ETC9
gives
$\Theta_b(y)=(b/y)\Theta_{1/b}(1/y)$.
As $y\to\infty$, $\Theta_b(y)\to1$ and
$y\Theta_b'(y),y^2\Theta_b'(y)\to0$.
As $y\downarrow0$, its leading term is $b/y$, with an
exponentially small remainder and differentiated remainders.
Consequently
\begin{equation}\tag{ETC15}
 \begin{gathered}
 \lim_{y\downarrow0}y^2\Theta_b'(y)=-b,\\
 \lim_{y\downarrow0}\bigl(y\Theta_b'(y)+\Theta_b(y)\bigr)=0,\\
 M_b(1)=\frac14[y^2\Theta_b'(y)]_{0}^{\infty}
       =\frac b4=\frac{a}{4N(k)},\\
 M_b(0)=\frac14[y\Theta_b'(y)+\Theta_b(y)]_{0}^{\infty}
       =\frac14.
 \end{gathered}
\end{equation}
Integrate first on a compact interval $[\epsilon,R]$ and then
take the endpoints; ETC9--10 guarantee convergence of the integrals,
and the displayed limits justify both passages. Thus the two endpoint
values are calculated from the actual reduced test, without assuming
special zeta values and without exchanging a nonabsolutely convergent
sum and integral.

The failure of the tempting degree interchange is quantified exactly.
Each $n\ne0$ term has integral zero in $dy$, by ETC4, but
\begin{equation}\tag{ETC16}
 \sum_{n\ne0}\int_0^\infty|\eta(ny/b)|dy
 =2b\left(\int_0^\infty|\eta(x)|dx\right)
                         \sum_{n=1}^\infty\frac1n=\infty.
\end{equation}
The coefficient integral is strictly positive because $\eta$ is
continuous and not identically zero. The codegree interchange also
lacks absolute convergence: each absolute integral with $dy/y$
equals the same positive finite number
$\int_0^\infty|\eta(x)|dx/x$, so its sum diverges. This does not
conflict with the convergent whole-test endpoint values ETC15.

\subsection{Sharp, conjugation, and two independent endpoint adjustments}

Use the source's literal involution
\begin{equation}\tag{ETC17}
 f^\sharp(u,y)=y^{-1}\overline{f(u^{-1},y^{-1})}.
\end{equation}
The tests $f_b$ are real and independent of $u$. Applying ETC9 at
$1/y$, and comparing with ETC8, proves
\begin{equation}\tag{ETC18}
 f_b^\sharp(y)=b f_{1/b}(y)
     =\mathcal E\widehat\xi_{k,a}(y),\quad
 M_{f_b^\sharp}(s)=bM_{1/b}(s)=M_b(1-s).
\end{equation}
For a general complex linear combination $f$ of these tests,
the exact formula instead is
\begin{equation}\tag{ETC19}
 M_{f^\sharp}(s)=\overline{M_f(1-\overline s)},\qquad
 M_{f^\sharp}(1)=\overline{M_f(0)}.
\end{equation}
Indeed substitute $v=1/y$ in the absolutely convergent integral
of $y^{-1}\overline{f(1/y)}y^{s-1}$; its integrand becomes
$\overline{f(v)}v^{-s}$, which is the displayed conjugated
Mellin transform. The source writes codegree both as $M_f(0)$
and as the degree of $f^\sharp$. Those values agree on its real
test functions, including every individual $f_b$ here; for complex
coefficients the latter is the conjugate of the former. The formulas
below specify which convention is being adjusted.

Define the complex-linear endpoint pair as
$(d(f),c(f))=(M_f(1),M_f(0))$. For two positive widths
$b_1\ne b_2$, both in the actual original family, ETC15 gives
\begin{equation}\tag{ETC20}
 \begin{pmatrix}d(c_1f_{b_1}+c_2f_{b_2})\\
                 c(c_1f_{b_1}+c_2f_{b_2})\end{pmatrix}
 =\frac14\begin{pmatrix}b_1&b_2\\1&1\end{pmatrix}
             \begin{pmatrix}c_1\\c_2\end{pmatrix},\qquad
 \det\left(\frac14\begin{pmatrix}b_1&b_2\\1&1\end{pmatrix}\right)
 =\frac{b_1-b_2}{16}\ne0.
\end{equation}
Thus prescribed changes $(D,C)\in\mathbb C^2$ are attained by
the explicit inverse
\begin{equation}\tag{ETC21}
 c_1=\frac{4(D-b_2C)}{b_1-b_2},\qquad
 c_2=\frac{4(b_1C-D)}{b_1-b_2}.
\end{equation}
Substitution gives $b_1c_1+b_2c_2=4D$ and $c_1+c_2=4C$,
proving the inverse without a dimension argument. Both widths may
use the same arbitrary divisor $k$ by choosing $a_i=N(k)b_i$.
The actual antecedent
$c_1\xi_{k,a_1}+c_2\xi_{k,a_2}$ still has zero value and zero
additive integral by ETC4, and its reduction is the displayed sum.
For real desired changes the coefficients are real. If instead
``codegree'' literally means $d(f^\sharp)$ for complex $f$,
replace $C$ by $\overline C$ in ETC21; ETC19 then proves the
corresponding real-linear inverse. This records the conjugation
rather than erasing it.

Finally the original adelic rational action is
\begin{equation}\tag{ETC22}
 (\alpha_r\xi)(x)=\xi(r^{-1}x),\quad
 \alpha_r\xi_{k,a}=\xi_{k+d(r),|r|_\infty a},\quad
 N(k+d(r))=|r|_\infty N(k),\quad
 \mathcal E\alpha_r\xi_{k,a}=\mathcal E\xi_{k,a}.
\end{equation}
The finite balls shift by their exact valuations; evenness of $\eta$
removes the real sign and leaves the real factor $|r|_\infty$.
The equality of $N$ follows from unique factorization of a nonzero
rational number. The last equality follows either from the retained
ratio in ETC6 or by reindexing the full rational sum $q\mapsto q/r$.
Real scaling alone changes $a$ to $\lambda a$, and hence
$f_b(y)$ to $f_{\lambda b}(y)=f_b(y/\lambda)$: its Mellin transform
acquires exactly $\lambda^s$, its degree exactly $\lambda$, and its
codegree remains unchanged. These are the actual two endpoint
coordinates used in ETC20--21, with every prime and real factor
retained. The calculation concerns this image family and its exact
adjustment maps; it does not assert positivity of a Weil pairing.

There is a two-dimensional adjustment space stable under the literal
sharp involution. Choose any $b>0$ with $b\ne1$ and put $c=1/b$.
In the basis $(f_b,f_c)$ the conjugate-linear operation is exactly
\begin{equation}\tag{ETC23}
 (z_1f_b+z_2f_c)^\sharp
 =\frac{\overline{z_2}}b f_b
       +b\overline{z_1}f_c,\qquad
 \begin{pmatrix}0&1/b\\b&0\end{pmatrix}^{\!2}=I.
\end{equation}
Both coefficients come from ETC18, including its width factor.
Define in this same space the original linear combinations
\begin{equation}\tag{ETC24}
 A_d=\frac{4(f_b-f_c)}{b-c},\qquad
 A_c=\frac{4(bf_c-cf_b)}{b-c}.
\end{equation}
Direct substitution into ETC15 proves
$(d(A_d),c(A_d))=(1,0)$ and $(d(A_c),c(A_c))=(0,1)$.
The determinant in ETC20 proves that these moments uniquely determine
an element of this two-dimensional space. ETC19 and its invariance
under sharp therefore give $A_d^\sharp=A_c$ and
$A_c^\sharp=A_d$, or equivalently
\begin{equation}\tag{ETC25}
 (D A_d+C A_c)^\sharp
       =\overline C A_d+\overline D A_c.
\end{equation}
These are explicit combinations of the two original tests and their
actual rational reductions. In particular they supply every prescribed
endpoint pair while retaining the involution and the original
zero-additive-moment antecedents.

For completeness, the two Gamma contributions in ETC11 have the
following separate behavior at the two endpoint arguments. Write
$h=s-1$ and $\zeta(s)=h^{-1}+c_0+O(h)$, where $c_0$ denotes its
retained Laurent constant. Set $L_b=\log b-\tfrac12\log\pi$.
The first Gamma summand including its factors has expansion
\begin{equation}\tag{ETC26}
 \begin{split}
 b^s\pi^{-s/2}\Gamma(s/2+2)\zeta(s)
 &=\frac{3b}{4h}+\frac{3b}{4}(c_0+L_b)
                     +\frac{b}{2\sqrt\pi}\Gamma'(5/2)+O(h),\\
 -\frac32b^s\pi^{-s/2}\Gamma(s/2+1)\zeta(s)
 &=-\frac{3b}{4h}-\frac{3b}{4}(c_0+L_b)
                     -\frac{3b}{4\sqrt\pi}\Gamma'(3/2)+O(h).
 \end{split}
\end{equation}
The Gamma recurrence differentiated at $3/2$ gives
$\Gamma'(5/2)=\tfrac32\Gamma'(3/2)+\Gamma(3/2)$.
Adding the two displayed expansions therefore retains and cancels
their opposite residues and their equal Laurent and width constants,
leaving exactly $b/4$. At $s=0$ the two summands have values
$\zeta(0)$ and $-\tfrac32\zeta(0)$. Their bracket is $-1/2$;
ETC15 and ETC11 imply $\zeta(0)=-1/2$, so their separate values
are $-1/2$ and $3/4$. Their sum is the independently calculated
$1/4$. Thus neither endpoint calculation has discarded the
contribution of one original Gaussian term.
